\subsection{Производящие функции: когда применять}
\begin{itemize}[leftmargin=*, nosep]
\item Здесь \textbf{обычные ПФ (OGF)}: $A(x)=\sum_{n\ge0}a_nx^n$, $[x^n]A=a_n$; коэффициенты с отрицательным индексом считаем нулевыми. Ряды формальные: аналитическая сходимость не нужна. В EGF коэффициент при $x^n$ равен $a_n/n!$ -- это другая нормировка.
\item Суммарный размер независимых выборов $\to$ произведение ПФ; упорядоченная последовательность непустых блоков $\to 1/(1-B)$; рекурсивное разбиение структуры $\to$ уравнение на ПФ.
\item Префиксы $\to A/(1-x)$; веса $n^k$ у коэффициентов $\to D^k A$, где $D=x\,d/dx$; линейная рекуррента $\to P/Q$; конечный автомат $\to (I-xT)^{-1}$.
\end{itemize}

\subsection{Операции и условия применимости}
\begin{itemize}[leftmargin=*, nosep]
\item Сложение: $[x^n](A+B)=a_n+b_n$. Произведение: $[x^n]AB=\sum_{k=0}^n a_kb_{n-k}$ -- \textbf{свёртка}, не $a_nb_n$.
\item Сдвиг: $[x^n]x^sA=a_{n-s}$ ($s\ge0$); удалить первые $s$ членов: $(A-\sum_{i<s}a_ix^i)/x^s$.
\item $A(rx)=\sum r^na_nx^n$; $A(x^d)$ оставляет коэффициенты только на кратных $d$ местах ($d\ge1$).
\item $A'=\sum_{n\ge0}(n+1)a_{n+1}x^n$; $DA=\sum na_nx^n$; $D^kA=\sum n^ka_nx^n$; $x^kA^{(k)}=\sum n^{\underline k}a_nx^n$.
\item $[x^n]A/(1-x)=\sum_{i=0}^na_i$; $[x^n]xA/(1-x)=\sum_{i=0}^{n-1}a_i$. Обратно, $(1-x)A$ кодирует $a_n-a_{n-1}$, $a_{-1}=0$.
\item $\int_0^x A(t)\,dt=\sum a_nx^{n+1}/(n+1)$; $a_n=A^{(n)}(0)/n!$ -- только когда соответствующие знаменатели обратимы.
\item $1/A$ существует, если $a_0$ обратим (над полем: $a_0\ne0$). Для композиции бесконечных рядов $A(B(x))$ достаточно $B(0)=0$; для полинома $A$ это ограничение не нужно. Композиционный обратный $F(G)=x$: $F(0)=0$, линейный коэффициент обратим.
\item Формулы с $\log$, $\exp$, дробными степенями понимаются над $\mathbb Q$ / полем характеристики 0. По модулю простого $p$ делить можно только на ненулевой остаток; для ряда до степени $N$ деления на $1,\ldots,N$ безопасны при $N<p$. Для $n!$, $n+1$ и других знаменателей условия проверяются отдельно.
\end{itemize}

\subsection{Базовые ПФ: строительные блоки}
Все последовательности индексируются с $n=0$, если не указано иначе; $k\ge0$, $m\ge1$.
\par\smallskip
{\renewcommand{\arraystretch}{1.65}\setlength{\tabcolsep}{3pt}
\begin{tabular}{@{}p{0.43\linewidth}p{0.52\linewidth}@{}}
\textbf{Коэффициент $a_n$} & \textbf{$\sum_{n\ge0}a_nx^n$}\\\hline
$[n=0]$; $[n=s]$ & $1$; $x^s$\\
$1$; $r^n$; $(-1)^n$ & $\frac1{1-x}$; $\frac1{1-rx}$; $\frac1{1+x}$\\
$[0\le n\le c]$ & $\frac{1-x^{c+1}}{1-x}$\\
$[n\equiv s\pmod d]$, $0\le s<d$ & $\frac{x^s}{1-x^d}$\\
$n$; $n^2$ & $\frac{x}{(1-x)^2}$; $\frac{x(1+x)}{(1-x)^3}$\\
$n^3$ & $\frac{x(1+4x+x^2)}{(1-x)^4}$\\
$n^4$ & $\frac{x(1+11x+11x^2+x^3)}{(1-x)^5}$\\
$\binom nk$; $n^{\underline k}$ & $\frac{x^k}{(1-x)^{k+1}}$; $\frac{k!x^k}{(1-x)^{k+1}}$\\
$\binom{n+m-1}{m-1}r^n$ & $\frac1{(1-rx)^m}$\\
$\binom mn r^n$, $0\le n\le m$ & $(1+rx)^m$\\
$1/n!$ & $e^x$\\
$1/n$ при $n\ge1$, $a_0=0$ & $-\log(1-x)$\\
$1/(n+1)$ & $-\frac{\log(1-x)}x$\\
$H_n=\sum_{i=1}^n1/i$, $H_0=0$ & $-\frac{\log(1-x)}{1-x}$\\
\end{tabular}}
\begin{itemize}[leftmargin=*, nosep]
\item Главная формула извлечения, $m\ge1$:
\[
[x^n]\frac{x^s}{(1-rx)^m}=
\begin{cases}\binom{n-s+m-1}{m-1}r^{n-s},&n\ge s,\\0,&n<s.\end{cases}
\]
\item Обобщённый бином: $(1+z)^\alpha=\sum_{n\ge0}\binom\alpha n z^n$, где $\binom\alpha n=\alpha^{\underline n}/n!$. В частности, $\binom{-m}n=(-1)^n\binom{n+m-1}{m-1}$. При подстановке ряда вместо $z$ его свободный член должен быть нулём.
\item Для полинома $p(n)$ степени $d$ ряд $\sum p(n)x^n$ рационален со знаменателем, делящим $(1-x)^{d+1}$. Получать через $D^k(1-x)^{-1}$ или разложение $p(n)$ по $\binom nk$.
\end{itemize}

\subsection{Базовые ПФ: комбинаторные последовательности}
{\renewcommand{\arraystretch}{1.75}\setlength{\tabcolsep}{3pt}
\begin{tabular}{@{}p{0.43\linewidth}p{0.52\linewidth}@{}}
\textbf{Последовательность} & \textbf{OGF / определяющее уравнение}\\\hline
Фибоначчи $F_0=0,F_1=1$ & $\frac{x}{1-x-x^2}$\\
$F_{n+1}$: композиции из 1 и 2 & $\frac1{1-x-x^2}$\\
Люка $L_0=2,L_1=1$ & $\frac{2-x}{1-x-x^2}$\\
Центральные биномы $\binom{2n}n$ & $\frac1{\sqrt{1-4x}}$\\
Каталан $C_n=\frac1{n+1}\binom{2n}n$ & $C=\frac{1-\sqrt{1-4x}}{2x}$; $C=1+xC^2$\\
Фусс--Каталан, $d\ge2$ & $T=1+xT^d$; $[x^n]T=\frac1{(d-1)n+1}\binom{dn}n$\\
Моцкин $1,1,2,4,9,21,\ldots$ & $M=1+xM+x^2M^2$\\
Большие Шрёдера $1,2,6,22,90,\ldots$ & $R=1+xR+xR^2$; $R=\frac{1-x-\sqrt{1-6x+x^2}}{2x}$\\
Композиции в положительные части: $a_0=1$, $a_n=2^{n-1}$ ($n\ge1$) & $\frac{1-x}{1-2x}$\\
Разбиения числа $p(n)$ & $\prod_{j\ge1}(1-x^j)^{-1}$\\
Разбиения на различные части & $\prod_{j\ge1}(1+x^j)$\\
Части $\le k$ (или не более $k$ частей) & $\prod_{j=1}^k(1-x^j)^{-1}$\\
Ровно $k$ положительных частей, без порядка & $x^k\prod_{j=1}^k(1-x^j)^{-1}$\\
Стирлинг II: $S(n,k)$, фиксированное $k\ge1$ & $\frac{x^k}{\prod_{j=1}^k(1-jx)}$\\
\end{tabular}}
\begin{itemize}[leftmargin=*, nosep]
\item В алгебраических уравнениях выбирается ряд со свободным членом 1. Для Моцкина: $M=(1-x-\sqrt{1-2x-3x^2})/(2x^2)$. Размер -- длина пути; шаги вверх/вниз имеют размер 1, горизонтальный -- тоже 1.
\item Частное с $x$ в знаменателе допустимо, если числитель делится на $x$ как ряд. Квадратный корень здесь берётся со свободным членом 1.
\item Для коэффициентов до $x^N$ бесконечные произведения по положительным размерам обрезаются на $j\le N$. Ряд различных частей также равен $\prod_{j\ge1}(1-x^{2j-1})^{-1}$ (разбиения на нечётные части).
\end{itemize}

\subsection{Независимые выборы, ограничения, порядок}
\begin{itemize}[leftmargin=*, nosep]
\item Одна категория веса $w>0$: выбрать 0/1 раз -- $1+x^w$; от 0 до $c$ раз -- $(1-x^{(c+1)w})/(1-x^w)$; любое число -- $(1-x^w)^{-1}$; хотя бы $a$ раз -- $x^{aw}/(1-x^w)$; число кратно $d$ -- $(1-x^{dw})^{-1}$.
\item Разрешённые размеры $S$ одного блока: $B(x)=\sum_{j\in S}b_jx^j$, где $b_j$ -- число вариантов размера $j$. Независимые категории перемножаются: размер суммируется, ответ -- нужный коэффициент произведения.
\item $k\ge1$ неотрицательных слагаемых с суммой $n$: $[x^n](1-x)^{-k}=\binom{n+k-1}{k-1}$; положительных: $[x^n]x^k(1-x)^{-k}=\binom{n-1}{k-1}$ при $n\ge k$, иначе 0.
\item $0\le a_i\le c_i$: $[x^n]\prod_i(1-x^{c_i+1})/(1-x)^k$. Нижние границы $a_i\ge l_i$ дают сдвиг на $\sum l_i$; с верхними границами заменить $c_i$ на $c_i-l_i$.
\item Подмножества с суммой весов: $\prod_i(1+x^{w_i})$; неограниченные количества типов: $\prod_i(1-x^{w_i})^{-1}$. Это выбор количеств \textbf{без порядка}, в отличие от последовательностей блоков.
\item Упорядоченные последовательности из ровно $k$ блоков -- $B^k$, из любого числа -- $1/(1-B)$, непустые -- $B/(1-B)$. Для бесконечной суммы нужно $B(0)=0$, иначе нулевые блоки могут дать бесконечно много объектов одного размера.
\item Произведение считает \textbf{пары выборов}. Чтобы оно считало требуемые объекты, разложение должно быть однозначным. Пересечение множеств вариантов само по себе не запрет, но если один итоговый объект получается несколькими способами, возникает кратный счёт.
\item Дополнительный параметр: $\prod_i(1+yx^{w_i})$; $[x^ny^k]$ считает подмножества веса $n$ из $k$ элементов. Маркер $y$ позволяет фиксировать число блоков/выбранных предметов.
\end{itemize}

\subsection{Линейные рекурренты и рациональные ПФ}
\begin{itemize}[leftmargin=*, nosep]
\item Если $a_n=c_1a_{n-1}+\cdots+c_ka_{n-k}$ при $n\ge k$, то $A=P/Q$, $Q=1-c_1x-\cdots-c_kx^k$, $\deg P<k$. Числитель: $P=QA\pmod{x^k}$, то есть $[x^j]P=a_j-\sum_{i=1}^jc_ia_{j-i}$ для $j<k$.
\item Обратно, $A=P/Q$, $Q=1+q_1x+\cdots+q_kx^k$: $a_n=-\sum_{i=1}^kq_ia_{n-i}$ для $n\ge k$, $n>\deg P$. Рациональная ПФ равносильна линейной рекурренте с постоянными коэффициентами начиная с некоторого номера.
\item Неоднородный член $b_n$ переводится в его ПФ с учётом первых индексов. Постоянный член обычно добавляет множитель $1-x$ в знаменатель до сокращения. Полиномиальная неоднородность тоже даёт рациональную ПФ.
\item Пример границы метода над $\mathbb Q$: $a_n=n^{2026}$ имеет рациональную ПФ; $a_n=1/(n+1)^2$ -- нет. Это не утверждение о рядах в конечной характеристике.
\item Один член $a_n$: DP $O(nk)$; Kitamasa $O(k^2\log n)$; плотная матрица $O(k^3\log n)$; Bostan--Mori для $[x^n]P/Q$ -- $O(M(k)\log n)$, где $M(k)$ -- стоимость умножения полиномов степени $O(k)$, $Q(0)$ обратим. Большую полиномиальную часть предварительно выделить.
\item Bostan--Mori: умножить числитель и знаменатель на $Q(-x)$; у $P(x)Q(-x)$ оставить коэффициенты чётности $n$, у $Q(x)Q(-x)$ -- чётные, сжать индексы вдвое, заменить $n$ на $\lfloor n/2\rfloor$. При $n=0$ ответ $P(0)/Q(0)$.
\end{itemize}

\subsection{Простые дроби и явная формула}
\begin{itemize}[leftmargin=*, nosep]
\item Если дробь неправильная, выделить $A=S+R/Q$, $\deg R<\deg Q$. Разложить $Q$; для каждого $(1-r_ix)^{m_i}$ нужны \textbf{все} степени $1,\ldots,m_i$:
\[
A=S+\sum_{i,j}\frac{c_{ij}}{(1-r_ix)^j},\qquad
 a_n=[x^n]S+\sum_{i,j}c_{ij}\binom{n+j-1}{j-1}r_i^n.
\]
\item Простой множитель даёт $r_i^n$; кратность $m_i$ -- полином по $n$ степени не выше $m_i-1$, умноженный на $r_i^n$ (в характеристике 0). Корень знаменателя -- $1/r_i$. Если множители не линейные над исходным полем, может понадобиться расширение поля; комплексная явная формула не обязательно удобна для вычисления по модулю.
\item Пример $a_n=2a_{n-1}+n$, $a_0=0$:
\[
 A=\frac{x}{(1-x)^2(1-2x)}
 =-\frac1{1-x}-\frac1{(1-x)^2}+\frac2{1-2x};\quad
 a_n=2^{n+1}-n-2.
\]
\end{itemize}

\subsection{Рекурсивные структуры: Каталан и Шрёдер}
\begin{itemize}[leftmargin=*, nosep]
\item ПСП единственно разбивается как \verb|(A)B|. Размер -- число пар скобок: $C=1+xC^2$, $C_0=1$, $C_n=\sum_{i=0}^{n-1}C_iC_{n-1-i}$. То же для упорядоченных бинарных деревьев с $n$ вершинами, где левый и правый потомки могут отсутствовать.
\item Решение $C=(1-\sqrt{1-4x})/(2x)$: знак выбирается по $C(0)=1$, другая ветвь не степенной ряд. $C_n=\frac1{n+1}\binom{2n}n=\binom{2n}n-\binom{2n}{n-1}$. Последняя форма не требует деления на $n+1$; сами биномы по модулю тоже нужно считать корректно.
\item Несколько независимых каталановских структур: для целого $r\ge1$, $[x^n]C(x)^r=\frac{r}{2n+r}\binom{2n+r}n$. Условия деления по модулю сохраняются.
\item В исходном конспекте Шрёдер сдвинут: $S_1=1,S_2=2,S_3=6$, $S(x)=\sum_{n\ge1}S_nx^n=xR(x)$. Поэтому $S=x+xS+S^2$ и $S=(1-x-\sqrt{1-6x+x^2})/2$, тогда как в таблице $R_0=1$.
\item При однозначном разложении на меньшие объекты: альтернативы складывать, независимые части перемножать, размер добавленных узлов отмечать степенью $x$. Пустой объект определяет свободный член.
\end{itemize}

\subsection{Примеры приёмов: блоки и суммы}
\begin{itemize}[leftmargin=*, nosep]
\item Двоичные строки с чётным числом максимальных блоков единиц. Ноль блоков разрешён: вклад $1/(1-x)$. При $2m$ блоках ($m\ge1$) есть $2m$ непустых серий единиц, $2m-1$ непустых внутренних серий нулей и две произвольные крайние серии:
\[
F=\frac1{1-x}+\sum_{m\ge1}\frac{x^{4m-1}}{(1-x)^{4m+1}}
=\frac{1-3x+3x^2}{1-4x+6x^2-4x^3}.
\]
\item Геометрическую сумму можно разложить как $F=\frac1{2(1-2x)}+\frac{1-2x}{2(1-2x+2x^2)}$. Коэффициенты с $n=0$: $1,1,1,2,6,16,36,72,\ldots$. Форма с $1/2$ не годится напрямую по модулю 2; первая рациональная форма годится. Альтернатива -- автомат по последнему биту и чётности числа блоков.
\item Производные и префиксы: $\sum_{n\ge0}n(n+1)x^n=2x/(1-x)^3$. Тогда
\[
\sum_{k=0}^n k(k+1)=[x^n]\frac{2x}{(1-x)^4}=\frac{n(n+1)(n+2)}3;
\]
\[
\sum_{k=1}^{n-1}k(k+1)=[x^n]\frac{2x^2}{(1-x)^4}=\frac{(n-1)n(n+1)}3.
\]
\item Верхняя граница $n-1$ вместо $n$ -- дополнительный множитель $x$ у ПФ префиксов. Сначала привести сумму к коэффициенту, затем применять формулу для $(1-x)^{-m}$.
\end{itemize}

\subsection{Автоматы и вычисление коэффициентов}
\begin{itemize}[leftmargin=*, nosep]
\item Пусть $v_n=Tv_{n-1}$, $v_0=b$ (векторы-столбцы, $T_{ij}$ -- число переходов $j\to i$). Тогда $V(x)=(I-xT)^{-1}b$. Для вектора принимающих состояний $c$: $A(x)=c^\top V(x)$; ПФ рациональна. Начальное состояние должно учитывать и пустую строку.
\item Для одного большого $n$ часто проще $c^\top T^n b$, чем символически обращать матрицу. Для фиксированного числа состояний можно перейти к скалярной рекурренте и применить методы выше.
\item Если нужны коэффициенты только до $x^N$, обычные суммы, произведения и обращение ряда с обратимым свободным членом можно считать по модулю $x^{N+1}$. При дифференцировании заранее нужна дополнительная точность: коэффициент $x^N$ производной зависит от $a_{N+1}$.
\item Свёртка через NTT требует подходящего модуля и длины преобразования (корней единицы нужного порядка). Для произвольного модуля -- подходящая свёртка через CRT/разбиение; обычное умножение работает без этих ограничений.
\end{itemize}
